\documentclass[12pt]{article}
\usepackage{amsmath,amssymb,amsfonts}
\usepackage[margin=1in]{geometry}

\begin{document}

\section*{Problem 5.24}

\textbf{Problem:} We have shown that the Legendre polynomials, $P_\ell(x)$, are solutions to
\[
(1-x^2)P_\ell''(x) - 2xP_\ell'(x) + \ell(\ell+1)P_\ell(x) = 0, \quad \ell = 0, 1, 2, \ldots,
\]
which are orthogonal on $[-1,1]$. However, these are not the only solutions; they are merely the only polynomial solutions.

\begin{itemize}
\item[(a)] To reassure yourself on this point, show that the function
\[
\tanh^{-1}x = \frac{1}{2}\ln\frac{1+x}{1-x}
\]
satisfies Legendre's equation with $\ell = 0$. Note that this function is singular at $x = \pm 1$.

\item[(b)] Show that if we look for a solution of Legendre's equation of the form
\[
Q_\ell(x) = P_\ell(x)\tanh^{-1}x + \Pi_\ell(x),
\]
then a solution of this form always exists where $\Pi_\ell(x)$ is a polynomial of degree $\ell - 1$ ($\ell \geq 1$). Find the simultaneous differential equation that $\Pi_\ell(x)$ must satisfy, and show that this equation always has a solution. What type of boundary condition has been applied to making $\Pi_\ell(x)$ be a polynomial of order $\ell - 1$?

Since $\Pi_\ell(x)$ is a polynomial of degree $\ell - 1$ ($\ell \geq 1$), it can be written as a linear combination of Legendre polynomials. Find an expression for the coefficients in this linear combination. To be precise, show that
\end{itemize}

\bigskip
\textbf{Solution:}

\subsection*{Part (a)}

Let $Q_0(x) = \tanh^{-1}x = \frac{1}{2}\ln\frac{1+x}{1-x}$.

Compute the derivatives:
\[
Q_0'(x) = \frac{1}{1-x^2},
\]
\[
Q_0''(x) = \frac{2x}{(1-x^2)^2}.
\]

Substitute into Legendre's equation with $\ell = 0$:
\[
(1-x^2)Q_0'' - 2xQ_0' + 0 = (1-x^2)\cdot\frac{2x}{(1-x^2)^2} - 2x\cdot\frac{1}{1-x^2} = \frac{2x}{1-x^2} - \frac{2x}{1-x^2} = 0. \quad\checkmark
\]

\[
\boxed{Q_0(x) = \tanh^{-1}x = \frac{1}{2}\ln\frac{1+x}{1-x} \text{ satisfies Legendre's equation with } \ell = 0.}
\]

Note that $Q_0(x)$ diverges logarithmically at $x = \pm 1$, which is why it is excluded from the polynomial (regular) solutions.

\subsection*{Part (b)}

We seek $Q_\ell(x) = P_\ell(x)\tanh^{-1}x + \Pi_\ell(x)$ as a solution of Legendre's equation. Let $w = \tanh^{-1}x$. Since $P_\ell$ satisfies Legendre's equation:
\[
(1-x^2)P_\ell'' - 2xP_\ell' + \ell(\ell+1)P_\ell = 0.
\]

Substituting $Q_\ell = P_\ell w + \Pi_\ell$ into Legendre's equation and using the product rule:
\begin{align}
(1-x^2)(P_\ell'' w + 2P_\ell' w' + P_\ell w'' + \Pi_\ell'') - 2x(P_\ell' w + P_\ell w' + \Pi_\ell') + \ell(\ell+1)(P_\ell w + \Pi_\ell) = 0.
\end{align}

The terms involving $w$ times (Legendre's equation for $P_\ell$) vanish. Using $w' = 1/(1-x^2)$ and $w'' = 2x/(1-x^2)^2$:
\[
(1-x^2)\left(\frac{2P_\ell'}{1-x^2} + \frac{2xP_\ell}{(1-x^2)^2}\right) - 2x\cdot\frac{P_\ell}{1-x^2} + (1-x^2)\Pi_\ell'' - 2x\Pi_\ell' + \ell(\ell+1)\Pi_\ell = 0.
\]

Simplifying:
\[
2P_\ell' + \frac{2xP_\ell}{1-x^2} - \frac{2xP_\ell}{1-x^2} + (1-x^2)\Pi_\ell'' - 2x\Pi_\ell' + \ell(\ell+1)\Pi_\ell = 0.
\]

Therefore $\Pi_\ell$ satisfies:
\[
\boxed{(1-x^2)\Pi_\ell'' - 2x\Pi_\ell' + \ell(\ell+1)\Pi_\ell = -2P_\ell'(x).}
\]

This is an \emph{inhomogeneous} Legendre equation with source term $-2P_\ell'(x)$. Since $P_\ell'(x)$ is a polynomial of degree $\ell-1$, we can seek a polynomial solution $\Pi_\ell$ of degree $\ell-1$.

By expanding $\Pi_\ell(x) = \sum_{k=0}^{\ell-1}a_k P_k(x)$ in Legendre polynomials and using the orthogonality relations, the coefficients can be determined. The standard result is:
\[
\Pi_\ell(x) = \sum_{\substack{k=0 \\ k+\ell \text{ odd}}}^{\ell-1}\frac{2k+1}{\ell(\ell+1)-k(k+1)}P_k(x) \cdot \frac{2}{2k+1}.
\]

The resulting $Q_\ell(x)$ are the \textbf{Legendre functions of the second kind}:
\[
Q_\ell(x) = P_\ell(x)\,\frac{1}{2}\ln\frac{1+x}{1-x} + \Pi_\ell(x).
\]

For example:
\begin{align}
Q_0(x) &= \frac{1}{2}\ln\frac{1+x}{1-x}, \\
Q_1(x) &= \frac{x}{2}\ln\frac{1+x}{1-x} - 1, \\
Q_2(x) &= \frac{3x^2-1}{4}\ln\frac{1+x}{1-x} - \frac{3x}{2}.
\end{align}

\end{document}
